\paragraph{Tableros residuales y su predictor~\cite[\S2 y \S6.1]{tc}.} El \emph{tablero ordinario} es el
tablero usual de \(n\times n\), vacío al comienzo. Un \emph{tablero residual} \(B\) es lo que deja de él
una colocación parcial de reinas sin ataques: conserva las coordenadas originales y consta de las \(m\)
filas vacías y las \(m\) columnas vacías. Una diagonal de cualquiera de las dos familias está
\emph{libre} si no contiene ninguna reina. El \emph{soporte} \(E(B)\) es el conjunto de las casillas que
están en una fila vacía y en una columna vacía y cuyas dos diagonales están libres; son las casillas
disponibles. Escribimos \(Q(B)\) para el número de soluciones que prolongan la colocación y \(B+e\) para
el hijo que resulta de poner una reina en \(e\in E(B)\). El tablero inicial \(B_n\) es el residual de la
colocación vacía, de modo que \(Q(B_n)=Q(n)\). Para pesos \(p=(p_e)_{e\in E(B)}\) y una diagonal libre
\(f\), sean \(\ell_f=\sum_{e\in f}p_e\) y \(h_f=1-\ell_f\). La entropía de Bethe de \(B\) es
\begin{equation}\label{eq:bethe-residual}
 H_B(p)=-\sum_e p_e\log p_e+3\sum_e(1-p_e)\log(1-p_e)-\sum_f h_f\log h_f
\end{equation}
en el dominio \(0<p_e<1/4\), \(h_f>0\), con masas de fila y de columna iguales a uno. Allí \(-H_B\) es
estrictamente convexa, de modo que \(H_B\) tiene a lo sumo un punto estacionario, el \emph{punto de
Bethe} de \(B\). Sea \(\mathcal A\) la matriz de incidencia de las casillas de \(E(B)\) con las líneas
de \(B\), a saber, sus filas vacías, sus columnas vacías y las diagonales libres que cortan a \(E(B)\),
y sea \(\tau_4(u)=u(1-u)/(1-4u)\). El operador de líneas es
\[
 S_B=\operatorname{diag}(0_R,0_C,h)+\mathcal A^{\mathsf T}\operatorname{diag}(\tau_4(p_e))\mathcal A .
\]
Cuando el grafo bipartito de filas y columnas del soporte es conexo, su núcleo está generado por
\((\mathbf1_R,-\mathbf1_C,0)\). El predictor de \(B\) es entonces
\begin{equation}\label{eq:predictor-residual}
 L(B)=H_B(p)-\tfrac12\sum_e\bigl[\log(1-4p_e)-4\log(1-p_e)\bigr]-\tfrac12\log\operatorname{pdet}S_B
      +\tfrac12\log(4\pi m),
\end{equation}
evaluado en el punto de Bethe, donde \(\operatorname{pdet}\) es el producto de los autovalores positivos.
Suprimir en \(S_B\) la fila y la columna de una misma línea axial da un cofactor igual a
\(\operatorname{pdet}S_B/(2m)\); usado con \(\tfrac12\log(2\pi)\), como en~\cite[\S\S6.1 y 7.1]{pc}, da
el mismo \(L(B)\). En el tablero inicial y para \(n\) grande, \(L(B_n)=\log Q_{\rm pred}(n)\), el
predictor del Apéndice~\ref{app:predictor}.

\paragraph{Historias y las dos leyes~\cite[\S2]{tc}.} Se sortea uniformemente un orden \(\sigma\) de las
filas. Bajo la ley uniforme \(\mu\) se elige una solución uniforme, de manera independiente de
\(\sigma\), y se revelan sus reinas en ese orden. Bajo la ley adaptativa \(\nu\), dada la siguiente fila
\(i\), de casillas \(R_i\), se elige la casilla \(e\in R_i\cap E(B)\) con probabilidad
\[
 \nu_B(e\mid i)=\frac{\exp L(B+e)}{\sum_{e'\in R_i\cap E(B)}\exp L(B+e')},\qquad
 \rho(B,i)=\sum_{e'\in R_i\cap E(B)}\exp\{L(B+e')-L(B)\}.
\]
Una \emph{historia} es la sucesión de tableros que resulta, junto con el orden de las filas. Un tablero
con \(m\) filas vacías tiene tamaño \(m\).

\paragraph{La regla de parada~\cite[def.~2.2]{tc}.} Se fijan \(0<\delta<1\), \(M=\lceil\delta n\rceil\) y
un entero \(K\) con \(K_0\le K\le M\). El umbral \(K_0\), las constantes y las tolerancias son los
de~\cite[\S5.4]{tc}. Sea \(B\) un tablero de una historia, de tamaño \(m\).
\begin{enumerate}
\item[(a)] \emph{Tramo macroscópico}, \(m\ge M\). El tablero es \emph{regular} si la discrepancia de la
historia al llegar a él, definida en~\cite[\S4.2]{tc} con la partición de paso \(a_{\rm mac}\)
de~\cite[prop.~4.11]{tc}, es menor que el radio \(\varrho_{\rm mac}\) de esa proposición.
\item[(b)] \emph{Corte}, \(m=M\). Además, los observables de~\cite[\S5.3]{tc} están dentro de sus
tolerancias de entrada y \(D_{\max}(M)<\alpha_{\rm hist}M\).
\item[(c)] \emph{Tramo mesoscópico}, \(K\le m<M\). El tablero es regular si su punto de Bethe existe y
pertenece a la clase interior: cumple las cotas de difusión, vacancia y
brecha~\eqref{eq:clase-importada} con las constantes fijadas en~\cite[\S5.4]{tc}, y la masa de cada
diagonal es a lo sumo \(\epsilon\).
\end{enumerate}
La \emph{primera salida} de una historia es su primer tablero no regular; su padre es regular. La
historia se detiene en la primera salida o, si no la hay, al llegar al tamaño \(K\). El suceso de
permanencia consiste en que no haya salida: todos los tableros hasta el de tamaño \(K\), incluido éste,
son regulares. En este artículo el suelo es \(K=j=\lfloor bn\rfloor\) (Sección~\ref{sec:signedmap}).
